\documentclass[11pt]{article}
\usepackage[margin=0.5in]{geometry}
\usepackage{tikz}
\usetikzlibrary{calc, decorations.markings, arrows.meta}

\begin{document}

% ============================================================
% Parametric Knit Stitch Geometry
% ============================================================
%
% A single knit stitch has this shape (viewed from front):
%
%          ___...---arc---...___
%         /                     \
%        / (tangent point)       \ (tangent point)
%       |                         |
%  left |   straight legs         | right
%  leg  |                         | leg
%       |                         |
%        \                       /
%         \___.---sweep---.___/   <-- lower arc, connects to next stitch
%
% The yarn path through a row visits each stitch left-to-right (or
% right-to-left on alternate rows), with the lower sweep connecting
% consecutive stitches.
%
% The legs of row N pass THROUGH the upper arc of the stitch in
% row N-1, creating the interlocking structure.

\begin{tikzpicture}[
  % === Controllable parameters ===
  % Horizontal spacing between stitch centers
  declare function={stitch_width=1.2;},
  % Half-width of the ellipse at the top of the stitch
  declare function={arc_rx=0.45;},
  % Half-height of the ellipse at the top of the stitch
  declare function={arc_ry=0.2;},
  % Length of the straight leg from base to tangent point
  declare function={leg_length=0.6;},
  % Gap between bottom of leg and top of lower arc
  declare function={leg_base_gap=0.05;},
  % Half-width of the lower connecting sweep
  declare function={sweep_rx=0.55;},
  % Half-height of the lower connecting sweep
  declare function={sweep_ry=0.15;},
  % Vertical spacing between rows (center to center)
  declare function={row_height=1.4;},
  % Yarn thickness
  yarn/.style={line width=1.2pt, draw=orange!80!black},
  yarn over/.style={yarn, draw=orange!80!black},
  yarn under/.style={yarn, draw=orange!80!black, opacity=0.3},
  % For crossing gaps (white overprint to show under-strand)
  crossing gap/.style={line width=3pt, draw=white},
]

% ---- Helper: draw a single knit stitch at position (#1, #2) ----
% #1 = x center of stitch
% #2 = y center (where the upper arc sits)
%
% The stitch consists of:
%   - Upper arc: top half of ellipse centered at (#1, #2)
%   - Left leg: from tangent point down to (#1 - arc_rx, #2 - leg_length)
%   - Right leg: from tangent point down to (#1 + arc_rx, #2 - leg_length)

% ==================================================================
% SHAPE 1: Cast-on tail
% ==================================================================
% A free end trailing from the first stitch
\node[anchor=east, font=\small\bfseries] at (-1.5, 0) {1. Cast-on tail};
\draw[yarn] (-0.8, -0.3) -- (0, 0);

% ==================================================================
% SHAPE 2: Cast-on loop (initial stitch on needle)
% ==================================================================
% A simple loop -- the seed for the first row
\node[anchor=east, font=\small\bfseries] at (-1.5, -1.8) {2. Cast-on loop};
\begin{scope}[shift={(0, -1.8)}]
  % Upper arc
  \draw[yarn] (-0.45, 0)
    arc[start angle=180, end angle=0, x radius=0.45, y radius=0.2];
  % Left leg down
  \draw[yarn] (-0.45, 0) -- (-0.45, -0.6);
  % Right leg down
  \draw[yarn] (0.45, 0) -- (0.45, -0.6);
\end{scope}

% ==================================================================
% SHAPE 3: Knit stitch body (the core shape)
% ==================================================================
% Two straight legs tangent to an ellipse, with the upper arc
\node[anchor=east, font=\small\bfseries] at (-1.5, -3.6) {3. Knit stitch};
\begin{scope}[shift={(0, -3.6)}]
  % Upper ellipse arc (top half)
  \draw[yarn] (-0.45, 0)
    arc[start angle=180, end angle=0, x radius=0.45, y radius=0.2];
  % Left leg: straight line from tangent point downward
  \draw[yarn] (-0.45, 0) -- (-0.45, -0.6);
  % Right leg: straight line from tangent point downward
  \draw[yarn] (0.45, 0) -- (0.45, -0.6);

  % Annotation: dimensions
  \draw[<->, thin, gray] (-0.45, 0.35) -- (0.45, 0.35)
    node[midway, above, font=\tiny] {$2 r_x$};
  \draw[<->, thin, gray] (0.7, 0) -- (0.7, -0.6)
    node[midway, right, font=\tiny] {$\ell$};
  \draw[<->, thin, gray] (-0.15, 0) -- (0.15, 0.2)
    node[midway, above left, font=\tiny] {$r_y$};
\end{scope}

% ==================================================================
% SHAPE 4: Inter-stitch sweep (lower arc connecting stitches)
% ==================================================================
% After the right leg descends, the yarn swoops to become
% the left leg of the next stitch via a lower elliptical arc
\node[anchor=east, font=\small\bfseries] at (-1.5, -5.4) {4. Inter-stitch};
\begin{scope}[shift={(0, -5.4)}]
  % First stitch right leg comes down to here
  \fill[black] (0.45, 0) circle(1.5pt);
  % Lower sweep arc connecting to next stitch's left leg
  \draw[yarn] (0.45, 0)
    arc[start angle=180, end angle=0, x radius=0.15, y radius=0.15];
  % Next stitch's left leg starts here
  \fill[black] (0.75, 0) circle(1.5pt);

  % Labels
  \node[font=\tiny, below] at (0.45, -0.2) {R leg};
  \node[font=\tiny, below] at (0.75, -0.2) {L leg};
\end{scope}

% ==================================================================
% SHAPE 5: Row-end turn (Bezier U-turn)
% ==================================================================
% At the end of a row, yarn wraps around and starts next row
\node[anchor=east, font=\small\bfseries] at (-1.5, -7.0) {5. Row-end turn};
\begin{scope}[shift={(0, -7.0)}]
  % Last stitch of row ends here
  \fill[black] (0, 0) circle(1.5pt);
  % Bezier curve wrapping around
  \draw[yarn] (0, 0)
    .. controls (0.8, 0) and (0.8, -1.4) .. (0, -1.4);
  % First stitch of next row starts here
  \fill[black] (0, -1.4) circle(1.5pt);

  \node[font=\tiny, right] at (0.85, -0.7) {selvage};
\end{scope}

% ==================================================================
% SHAPE 6: Interlocking (new row through old row)
% ==================================================================
% The legs of the new stitch pass through the old stitch's loop
\node[anchor=east, font=\small\bfseries] at (-1.5, -9.5) {6. Interlocking};
\begin{scope}[shift={(0, -9.5)}]
  % OLD stitch (row below) -- upper arc
  \draw[yarn, blue!60!black] (-0.45, 0)
    arc[start angle=180, end angle=0, x radius=0.45, y radius=0.2];
  % Old stitch legs
  \draw[yarn, blue!60!black] (-0.45, 0) -- (-0.45, -0.6);
  \draw[yarn, blue!60!black] (0.45, 0) -- (0.45, -0.6);

  % NEW stitch (row above) -- its legs pass THROUGH the old loop
  % Left leg of new stitch (comes from above, goes under old arc)
  % Draw crossing gap on old stitch where new leg crosses it

  % New stitch upper arc (sits above old stitch)
  \draw[yarn, red!60!black] (-0.35, 0.9)
    arc[start angle=180, end angle=0, x radius=0.35, y radius=0.18];

  % New left leg: above crossing, then gap, then below
  % Over segment (above old arc)
  \draw[yarn, red!60!black] (-0.35, 0.9) -- (-0.35, 0.25);
  % Gap in old stitch arc (white overprint)
  \draw[crossing gap] (-0.35, 0.12) -- (-0.35, -0.05);
  % Under segment (below old arc, going to bottom)
  \draw[yarn, red!60!black] (-0.35, 0.12) -- (-0.35, -0.4);

  % New right leg
  \draw[yarn, red!60!black] (0.35, 0.9) -- (0.35, 0.25);
  \draw[crossing gap] (0.35, 0.12) -- (0.35, -0.05);
  \draw[yarn, red!60!black] (0.35, 0.12) -- (0.35, -0.4);

  % Labels
  \node[font=\tiny, blue!60!black, right] at (0.5, -0.3) {old loop};
  \node[font=\tiny, red!60!black, right] at (0.5, 0.6) {new stitch};

  % Crossing markers
  \node[font=\tiny, fill=green!20, draw=green!60!black, circle, inner sep=1pt]
    at (-0.35, 0.1) {+};
  \node[font=\tiny, fill=green!20, draw=green!60!black, circle, inner sep=1pt]
    at (0.35, 0.1) {+};
\end{scope}

% ==================================================================
% SHAPE 7: Bind-off / final tail
% ==================================================================
\node[anchor=east, font=\small\bfseries] at (-1.5, -11.5) {7. Final tail};
\begin{scope}[shift={(0, -11.5)}]
  % Last loop
  \draw[yarn] (-0.45, 0)
    arc[start angle=180, end angle=0, x radius=0.45, y radius=0.2];
  \draw[yarn] (-0.45, 0) -- (-0.45, -0.4);
  \draw[yarn] (0.45, 0) -- (0.45, -0.4);
  % Tail going through last loop and trailing off
  \draw[yarn, dashed] (0, 0.2) -- (0, 0.8)
    node[above, font=\tiny] {through last loop};
\end{scope}

% ==================================================================
% COMBINED: A complete row of 3 knit stitches
% ==================================================================
\begin{scope}[shift={(5, 0)}]
  \node[anchor=south west, font=\small\bfseries] at (-0.5, 1.2)
    {Combined: 3 stitches with cast-on tail};

  % Cast-on tail
  \draw[yarn] (-1.0, -0.6) -- (-0.45, -0.6);

  % === Stitch 0 (x=0) ===
  % Upper arc
  \draw[yarn] (-0.45, 0)
    arc[start angle=180, end angle=0, x radius=0.45, y radius=0.2];
  % Legs
  \draw[yarn] (-0.45, 0) -- (-0.45, -0.6);
  \draw[yarn] (0.45, 0) -- (0.45, -0.6);

  % Inter-stitch sweep 0->1
  \draw[yarn] (0.45, -0.6)
    arc[start angle=180, end angle=0, x radius=0.15, y radius=0.15];

  % === Stitch 1 (x=1.2) ===
  \draw[yarn] ({1.2-0.45}, 0)
    arc[start angle=180, end angle=0, x radius=0.45, y radius=0.2];
  \draw[yarn] ({1.2-0.45}, 0) -- ({1.2-0.45}, -0.6);
  \draw[yarn] ({1.2+0.45}, 0) -- ({1.2+0.45}, -0.6);

  % Inter-stitch sweep 1->2
  \draw[yarn] ({1.2+0.45}, -0.6)
    arc[start angle=180, end angle=0, x radius=0.15, y radius=0.15];

  % === Stitch 2 (x=2.4) ===
  \draw[yarn] ({2.4-0.45}, 0)
    arc[start angle=180, end angle=0, x radius=0.45, y radius=0.2];
  \draw[yarn] ({2.4-0.45}, 0) -- ({2.4-0.45}, -0.6);
  \draw[yarn] ({2.4+0.45}, 0) -- ({2.4+0.45}, -0.6);

  % Row-end turn
  \draw[yarn] ({2.4+0.45}, -0.6)
    .. controls ({2.4+1.0}, -0.6) and ({2.4+1.0}, -2.0) ..
    ({2.4+0.45}, -2.0);

  % === Row 2: stitches 3,4,5 going right to left ===
  % These interlock with row 1 above them

  % Stitch 5 (x=2.4, row 2) -- first worked in row 2 (right to left)
  % Upper arc at y=-1.4 (sits within stitch 2's legs)
  \draw[yarn] ({2.4+0.35}, -1.4)
    arc[start angle=0, end angle=180, x radius=0.35, y radius=0.18];
  % Left leg passes through stitch 2's loop
  \draw[yarn] ({2.4-0.35}, -1.4) -- ({2.4-0.35}, -2.0);
  % Right leg passes through stitch 2's loop
  \draw[yarn] ({2.4+0.35}, -1.4) -- ({2.4+0.35}, -2.0);

  % Crossing gaps on stitch 2's legs where stitch 5's arc crosses
  % (stitch 5's upper arc passes BEHIND stitch 2's legs)

  % Inter-stitch sweep 5->4 (going right to left now)
  \draw[yarn] ({2.4-0.35}, -2.0)
    arc[start angle=0, end angle=180, x radius=0.175, y radius=0.15];

  % Stitch 4 (x=1.2, row 2)
  \draw[yarn] ({1.2+0.35}, -1.4)
    arc[start angle=0, end angle=180, x radius=0.35, y radius=0.18];
  \draw[yarn] ({1.2-0.35}, -1.4) -- ({1.2-0.35}, -2.0);
  \draw[yarn] ({1.2+0.35}, -1.4) -- ({1.2+0.35}, -2.0);

  % Inter-stitch sweep 4->3
  \draw[yarn] ({1.2-0.35}, -2.0)
    arc[start angle=0, end angle=180, x radius=0.175, y radius=0.15];

  % Stitch 3 (x=0, row 2)
  \draw[yarn] ({0+0.35}, -1.4)
    arc[start angle=0, end angle=180, x radius=0.35, y radius=0.18];
  \draw[yarn] ({0-0.35}, -1.4) -- ({0-0.35}, -2.0);
  \draw[yarn] ({0+0.35}, -1.4) -- ({0+0.35}, -2.0);

  % Row labels
  \node[font=\tiny, gray] at (-0.8, 0) {Row 1 $\rightarrow$};
  \node[font=\tiny, gray] at (-0.8, -1.4) {$\leftarrow$ Row 2};
\end{scope}

\end{tikzpicture}

\end{document}
